\subsection{向量函数的中值定理}

定理 2. 设 \(\mathbf{f}\) 是向量函数，定义在有界闭区间 \(\mathrm{I} = [a, b]\) （\(\mathbf{R}\) 中的；其中 \(a < b\) ）上且连续，取值于一个赋范空间 \(\mathrm{E}\) （\(\mathbf{R}\) 上的）中；设 \(g\) 是 \(\mathrm{I}\) 上的连续递增实函数。设 \(\mathbf{f}\) 与 \(g\) 在 \([a, b]\) 中该区间的一可数子集 A 的相对余集的各点处具有右导数（允许 \(g_r'(x)\) 在某些点 \(x \notin \mathbf{A}\) 处为无穷），并设在每一这样的点处有

\[
\left| \left| \mathbf {f} _ {d} ^ {\prime} (x) \right| \right| \leqslant g _ {r} ^ {\prime} (x).\tag{6}
\]

在这些假设下，有

\[
\| \mathbf {f} (b) - \mathbf {f} (a) \| \leqslant g (b) - g (a).\tag{7}
\]

证明与命题 2 的证明类似。任取 \(\varepsilon > 0\) ，并设 \((a_{n})\) 是把 A 按某种次序枚举而得的序列。设 J 是由这样的点 \(y \in I\) 所成的集：对所有的 \(x\) （\(a \leqslant x \leqslant y\) ）有

\[
\| \mathbf {f} (x) - \mathbf {f} (a) \| \leqslant g (x) - g (a) + \varepsilon (x - a) + \varepsilon \sum_ {a _ {n} <   x} \frac {1}{2 ^ {n}};\tag{8}
\]

我们将证明 J = I。与命题 2 中一样，立即可以看出 J 是以 a 为左端点的区间；若 c 是其右端点，则 \(c \in J\) ；事实上，对所有 x \textless{} c 有 (8)，从而更有

\[
\| \mathbf {f} (x) - \mathbf {f} (a) \| \leqslant g (c) - g (a) + \varepsilon (c - a) + \varepsilon \sum_ {a _ {n} <   c} \frac {1}{2 ^ {n}}
\]

由此，在此不等式中让 \(x\) 趋于 \(c\) ，由 \(\mathbf{f}\) 的连续性可知 \(c\) 满足 (8)。

我们来证明必有 c = b。设 c \textless{} b 且进而 \(c \notin A\) ：则 \(\mathbf{f}_{d}^{\prime}(c)\) 与 \(g_{r}^{\prime}(c)\) 存在并满足 (6)；首先假设 \(g_{r}^{\prime}(c)\) （它必然 \(\geqslant 0\) ）为有限；则总可写 \(\mathbf{f}_{d}^{\prime}(c) = \mathbf{u}g_{r}^{\prime}(c)\) ，其中 \(\|u\| \leqslant 1\) ；由于函数 \(\mathbf{f}(x) - \mathbf{u}g(x)\) 在点 c 的右导数为零，必存在 y 使 \(c < y \leqslant b\) ，且对 \(c \leqslant x \leqslant y\) 有

\[
\| \mathbf {f} (x) - \mathbf {f} (c) - \mathbf {u} (g (x) - g (c)) \| \leqslant \varepsilon (x - c)
\]

由此

\[
\| \mathbf {f} (x) - \mathbf {f} (c) \| \leqslant g (x) - g (c) + \varepsilon (x - c)
\]

并且，考虑到 (8) （其中 x 换为 c），

\[
\begin{array}{c} \| \mathbf {f} (x) - \mathbf {f} (a) \| \leqslant g (x) - g (a) + \varepsilon (x - a) + \varepsilon \sum_ {a _ {n} <   c} \frac {1}{2 ^ {n}} \\ \leqslant g (x) - g (a) + \varepsilon (x - a) + \varepsilon \sum_ {a _ {n} <   x} \frac {1}{2 ^ {n}}. \end{array}
\]

于是有 \(y \in J\) ，这是一个矛盾。其次设 \(c \notin A\) 且 \(g_{r}^{\prime}(c) = +\infty\) ；则存在 y 使 \(c < y \leqslant b\) ，且对 \(c \leqslant x \leqslant y\) 一方面有

\[
\left\| \mathbf {f} (x) - \mathbf {f} (c) \right\| \leqslant \left(\left\| \mathbf {f} _ {d} ^ {\prime} (c) \right\| + 1\right) (x - c)
\]

而另一方面

\[
g (x) - g (c) \geqslant \left(\left\| \mathbf {f} _ {d} ^ {\prime} (c) \right\| + 1\right) (x - c)
\]

由此

\[
\| \mathbf {f} (x) - \mathbf {f} (c) \| \leqslant g (x) - g (c)
\]

并且可如上一样得出结论。最后，若有 \(c = a_{k}\) ，则存在 y 使 \(c < y \leqslant b\) ，且对 \(c < x \leqslant y\) 有

\[
\| \mathbf {f} (x) - \mathbf {f} (c) \| \leqslant \frac {\varepsilon}{2 ^ {k}}
\]

由此，考虑到 (8) （其中 x 换为 c），

\[
\begin{array}{c} \| \mathbf {f} (x) - \mathbf {f} (a) \| \leqslant g (c) - g (a) + \varepsilon (c - a) + \varepsilon \sum_ {a _ {n} <   x} \frac {1}{2 ^ {n}} \\ \leqslant g (x) - g (a) + \varepsilon (x - a) + \varepsilon \sum_ {a _ {n} <   x} \frac {1}{2 ^ {n}} \end{array}
\]

这又导出一个矛盾。证明的结尾与命题 2 的相同。

证毕。

注。1) 这里同样地，在定理 2 的陈述中，可把区间 {[}a, b{[} 换成 {]}a, b{]} ，并把“右导数”换成“左导数”。

\begin{enumerate}
\def\labelenumi{\arabic{enumi})}
\setcounter{enumi}{1}
\tightlist
\item
  我们将在后面说明如何确定 (7) 中等号成立的情形，以及如何把定理 2 推广到 E 为任意局部凸空间的情形，这要借助另一种证明方法，它使人们能从定理 1 推出定理 2。
\end{enumerate}

推论. 为使区间 \(I \subset R\) 上的、取值于 R 上的赋范空间 E 中的连续向量函数在 I 上为常数，只需它在 I 的一可数子集的（相对于 I 的）余集的各点处右导数为零。

注. 定理 1 与定理 2 的证明本质上依赖于域 R 的特殊拓扑性质；可以给出一些赋值域 K 的例子，对它们存在从 K 到其自身的非常数线性映射，它在每一点的导数都为零（cf.~I, p.~37, exerc. 2）。

命题 3. 设 f 是取值于 R 上的赋范空间 E 中的向量函数，定义在区间 I ⊂ R 上且连续，并在 I 的一可数子集的（相对于 I 的）余集 B 上右可导；则对所有的点 \(x_{0} \in B\) 、\(x \in I\) 、\(y \in I\) ，有（例如设 x \textless{} y）

\[
\left\| \mathbf {f} (y) - \mathbf {f} (x) - \mathbf {f} _ {d} ^ {\prime} (x _ {0}) (y - x) \right\| \leqslant (y - x) \sup _ {z \in \mathrm{B}, x <   z <   y} \left\| \mathbf {f} _ {d} ^ {\prime} (z) - \mathbf {f} _ {d} ^ {\prime} (x _ {0}) \right\|.\tag{9}
\]

事实上，只需应用定理 2，把 f 换成函数

\[
\mathbf {f} (z) - \mathbf {f} _ {d} ^ {\prime} (x _ {0}) z,
\]

并把 \(g\) 换成导数为 \(\sup_{z\in \mathrm{B}, x < z < y}\left\| \mathbf{f}_d'(z) - \mathbf{f}_d'(x_0)\right\|\) 的线性函数。

定理 2 可推广到复变量的向量函数：

命题 4. 设 \(\mathbf{f}\) 是复变量的连续可导函数，定义在域 \(\mathbf{C}\) 的凸开子集 A 上，取值于域 \(\mathbf{C}\) 上的一个赋范空间 E 中。若有 \(\| \mathbf{f}'(z)\| \leqslant m\) （对所有 \(z\in \mathrm{A}\) ），则有 \(\| \mathbf{f}(b) - \mathbf{f}(a)\| \leqslant m|b - a|\) （对 A 的每一对点 \(a\) 、\(b\) ）。

我们令 \(\mathbf{g}(t) = \frac{1}{b - a}\) \(\mathbf{f}(a + t(b - a))\) （对 \(0 \leqslant t \leqslant 1\) ）；由于 \(\mathbf{g}'(t) = \mathbf{f}'(a + t(b - a))\) ，对函数 \(\mathbf{g}\) 应用定理 2 立即得到命题。

推论. 为使复变量的向量函数 f（定义并连续于开集 \(A \subset C\) 上、取值于 C 上的赋范空间中）为常数，只需它在 A 的每一点处导数为零。

事实上，设 a 是 A 的任意一点；A 中使 \(\mathbf{f}(z) = \mathbf{f}(a)\) 的点 z 所成的集 B 因 f 连续而是闭集；它也是开集，这可对 B 的任意一点的一个包含在 A 中的凸开邻域应用命题 4 （取 m = 0）看出；故它与 A 恒同。

命题 5. 设 f 是复变量的向量函数，定义、连续且可导于凸开集 A ⊂ C 上，取值于域 C 上的赋范空间中；则不论 A 中的点 \(x_{0}\) 、x 与 y 为何，有

\[
\left| \left| \mathbf {f} (y) - \mathbf {f} (x) - \mathbf {f} _ {d} ^ {\prime} (x _ {0}) (y - x) \right| \right| \leqslant | y - x | \sup _ {z \in A} \left| \left| \mathbf {f} ^ {\prime} (z) - \mathbf {f} ^ {\prime} (x _ {0}) \right| \right|.\tag{10}
\]

只需对函数

\[
\mathbf {g} (t) = \mathbf {f} (x + t (y - x)) - \mathbf {f} ^ {\prime} (x _ {0}) (y - x) t
\]

在区间 {[}0, 1{]} 上应用定理 2。

